\section{Background}
\label{sec:background}

\subsection{Offline actor--critic learning}

We consider a discounted Markov decision process
$\langle\mathcal S,\mathcal A,P,r,\gamma\rangle$ with continuous action space
$\mathcal A\subset\mathbb R^d$. Offline RL receives only a fixed dataset
$\mathcal D=\{(s_i,a_i,r_i,s_i')\}_{i=1}^N$ and must learn a policy without
additional data collection. A critic $\widehat Q$ is trained from bootstrapped
targets of the form
\begin{equation}
 y_{\rho}=r+\gamma\min_{j}Q^-_j(s',\rho(s')),
 \label{eq:td_target_background}
\end{equation}
where $\rho$ is the target policy and $Q^-_j$ are target critics. When $\rho$
selects actions outside dataset support, errors in $Q^-_j$ enter both the
regression target and the next actor update. This feedback is a central source
of offline instability \citep{fujimoto2019offpolicy,kumar2019bear}.

Behavior-anchored algorithms control this feedback by optimizing a value term
together with a cloning or divergence penalty. For a deterministic actor
$\mu_\theta$, the normalized TD3+BC form is
\begin{equation}
 \mathcal L_{\rm actor}(\theta)
 =-\alpha\,\mathbb E_s\bar Q(s,\mu_\theta(s))
 +\mathbb E_{(s,a)\sim\mathcal D}\frac{1}{d}
   \|\mu_\theta(s)-a\|_2^2,
 \label{eq:td3bc_background}
\end{equation}
where $\bar Q=\widehat Q/C$ and $C>0$ is a stop-gradient critic scale. The
anchor makes the update conservative, but a one-step method uses the same
policy branch for both roles: a Polyak target copy enters the TD target, while
the online actor is deployed. Consequently, the displacement that perturbs
bootstrapping is tied to the total displacement used for policy improvement.

\subsection{Proximal updates and Wasserstein geometry}

Fix a state distribution $\rho$ and a policy family $\mathcal M$. Given an
estimated energy $\widehat{\mathcal E}$ and a metric $d_{\mathcal M}$, a
proximal policy update is
\begin{equation}
 \pi^+\in\argmin_{\pi\in\mathcal M}
 \left\{\widehat{\mathcal E}(\pi)
 +\frac{1}{2h}d_{\mathcal M}^2(\pi,\pi^{\rm ref})\right\}.
 \label{eq:generic_prox}
\end{equation}
This is the minimizing-movement, or implicit-Euler, construction associated
with a metric gradient flow \citep{ambrosio2005gradient,otto2001geometry}. The
variational problem in Eq.~\eqref{eq:generic_prox} is meaningful for finite
$h$ whenever a minimizer exists. Its interpretation as an accurate time step
of a continuous flow is more restrictive: it requires a fixed energy,
sufficiently small realized moves, and increasingly accurate solves as the
step size vanishes.

For state-conditioned deterministic policies
$\pi_\mu(\cdot\mid s)=\delta_{\mu(s)}$, the quadratic Wasserstein distance has
the Dirac limit
\begin{equation}
 W_2^2\!\left(\pi_\mu(\cdot\mid s),\pi_\nu(\cdot\mid s)\right)
 =\|\mu(s)-\nu(s)\|_2^2.
 \label{eq:dirac_w2_background}
\end{equation}
Thus squared action anchoring is a Wasserstein movement cost, up to a fixed
metric scaling. We use this observation to separate three levels of analysis:
(i) a local Wasserstein-flow limit, (ii) finite-step inexact proximal
improvement relative to the learned critic, and (iii) the coupled actor--critic
branch dynamics that dominate at large policy-improvement budgets.
